\section{The Fermat torus cover}\label{sec:fer}

The base $\Gamma$ is the hexagonal torus $(\R/\Z)^2$, with $z=u+\zeta v$ and $\zeta=e^{2\pi i/3}$.
This section constructs the unit Fermat lift $Z(u,v)$, the two deck transformations of $\R^2\times\C$, the model form and
the handle data used in \S\ref{sec:cub}, and the first Moser primitive $\alpha_1$, which vanishes as an ambient covector
on the three marked fibres. The base area three is computed from the quasi-periods of the lift.
We write $\mathcal F(Z)=Z_0^3+Z_1^3+Z_2^3$, $S_\Gamma=\{|Z|=1,\ \mathcal F(Z)=0\}\subset\C^3$, and
$\eta_H=\operatorname{Im}(Z^\dagger dZ)/(2\pi)$ on the unit sphere.

\subsection{The representative and cover data}\label{sec:fer-curve}

\begin{lemma}[Unit Fermat representatives]\label{lem:fer-curve}
\lean{SymplecticPacking.Embed.Fpoly_Zunit,SymplecticPacking.Bundle.normSq3_Zunit,SymplecticPacking.Bundle.contDiff_Zunit}
The theta lift defined below is smooth, has unit norm and satisfies
$Z_0(u,v)^3+Z_1(u,v)^3+Z_2(u,v)^3=0$ for every $(u,v)\in\R^2$.
\end{lemma}

\begin{proof}
Away from the lattice the holomorphic lift $\hat Z$ is a nonzero scalar multiple of the Weierstrass vector of
Lemma~\ref{lem:fer-torus}, so it satisfies the homogeneous cubic equation.
At the lattice it is a nonzero multiple of $(0,1,-1)$.
Normalization preserves that equation, gives unit norm, and is smooth because $\hat Z$ never vanishes
(Lemma~\ref{lem:fer-theta}).
\end{proof}

We use the cover model of Definition~\ref{def:cub-model}: $\widetilde N=\R^2\times\C$ with the deck maps $T_u$, $T_v$, the
map $\mathsf N(u,v,w)=(Z(u,v),w)$, the gauge relation $(Z,U)\sim(e^{i\chi}Z,e^{3i\chi}U)$, $p=\pi|w|^2$,
$\eta=\operatorname{Im}\sum_j\bar Z_jdZ_j/(2\pi)$ and $\sigma=d\eta$; here $Z(u+1,v)=-Z(u,v)$, $Z(u,v+1)=\phi(u,v)Z(u,v)$ and
$|\phi|=1$.

\begin{lemma}[Angular and model forms]\label{lem:fer-alpha}
\lean{SymplecticPacking.Bundle.d_etaH,SymplecticPacking.Return.d1_betaLoc}
On $w\ne0$, the form $\alpha=d\arg w/(2\pi)-3\eta$ is deck invariant and $d\alpha=-3\sigma$.
It has integral one on each fibre circle. The product $p\alpha$ extends smoothly across $w=0$.
\end{lemma}

\begin{proof}
Under a variable phase $Z\mapsto e^{i\chi}Z$, $\eta$ changes by $d\chi/(2\pi)$.
The fibre coordinate changes by $e^{3i\chi}$, whose angular form changes by $3d\chi/(2\pi)$.
These changes cancel. The remaining claims follow from $d\eta=\sigma$ and
$p\,d\arg w/(2\pi)=\frac12\operatorname{Im}(\bar w\,dw)$.
\end{proof}

\begin{lemma}[The coordinate model form]\label{lem:fer-Omega0}
\lean{SymplecticPacking.Bundle.Omega0loc,SymplecticPacking.Return.d1_betaLoc,SymplecticPacking.Bundle.Cubic.OmegaT_split,SymplecticPacking.Bundle.Cubic.OmegaT_sq}
On the cover,
\begin{equation}\label{eq:fer-Omega0-local}
 \Omega_0=\sigma+d(p\alpha)
 =\omega_V+(1-3p)\sigma-3dp\wedge\eta.
\end{equation}
It is closed, deck invariant and smooth at the fibre axis. Its square is
$2(1-3p)\omega_V\wedge\sigma$, so it is nondegenerate for $p<1/3$.
\end{lemma}

\begin{proof}
Differentiate the smooth primitive
$\frac12\operatorname{Im}(\bar w\,dw)+(1-3p)\eta$.
In its square, every additional mixed term has a repeated vertical differential or three base differentials.
Positivity of $\sigma$ is Lemma~\ref{lem:fer-sigma}.
\end{proof}

\subsection{An explicit parametrisation and atlas}\label{sec:fer-torus}

We fix an explicit uniformisation of $\Gamma$ by the torus $\C/\Lambda_0$, $\Lambda_0=\Z+\zeta\Z$,
$\zeta=e^{2\pi i/3}$, with real coordinates $(u,v)\in(\R/\Z)^2$, $z=u+\zeta v$, and a lift of the curve to the unit sphere
without poles. In these coordinates $\sigma$, $\eta_H$ and all base data are explicit functions of $(u,v)$, smooth, with the flat cutoffs of \S\ref{sec:fer-handle}. Since $\zeta_6:=e^{i\pi/3}=1+\zeta$, also $\Lambda_0=\Z+\zeta_6\Z$, and
$\zeta_6\Lambda_0=\Lambda_0$ because $\zeta_6^2=\zeta_6-1$; likewise $\zeta\Lambda_0=\Lambda_0$.

\begin{definition}[Torus parametrisation]\label{def:fer-torus}
\lean{SymplecticPacking.Bundle.tau, SymplecticPacking.Bundle.hexPair,
SymplecticPacking.Bundle.mem_lattice_iff, SymplecticPacking.Bundle.torusZ,
SymplecticPacking.Bundle.cF, SymplecticPacking.Bundle.dF, SymplecticPacking.Bundle.torusLift,
SymplecticPacking.Bundle.latticeLift, SymplecticPacking.Bundle.torusPoint}
Let $\wp=\wp_{\Lambda_0}$ be the Weierstrass function of $\Lambda_0$, $w_\wp=\wp'(z)$, and $g_2,g_3$ the invariants of
$\Lambda_0$, so that $w_\wp^2=4\wp^3-g_2\wp-g_3$. Put
\[
 c_F=g_3^{-1/3},\qquad d_F=(3g_3)^{-1/2}\qquad\text{(principal powers)},
\]
so that $c_F^3g_3=1$ and $3d_F^2g_3=1$ ($g_3\ne0$ by Lemma~\ref{lem:fer-torus}). Define $\Phi_\Gamma$ on the quotient by the lattice, with brackets denoting nonzero complex vectors modulo a scalar by
\[
 \Phi_\Gamma(z)=\big[\,2c_F\,\wp(z):-(1+d_F\,w_\wp):-(1-d_F\,w_\wp)\,\big]\quad(z\notin\Lambda_0),\qquad
 \Phi_\Gamma(0)=[0:1:-1].
\]
The torus is oriented by $\nu=du\wedge dv$, so $\int\nu=1$.
\end{definition}

The orientation of $\nu$ is the complex one: $d\operatorname{Re}z\wedge d\operatorname{Im}z=\operatorname{Im}\zeta\,du\wedge dv$ and
$\operatorname{Im}\zeta=\sqrt3/2>0$. The constants $c_F,d_F$ normalise the map rather than the lattice.

\begin{lemma}[Injectivity of the Fermat representative map]\label{lem:fer-torus}
\lean{SymplecticPacking.Bundle.hex_g2, SymplecticPacking.Bundle.hex_g3_ne_zero,
SymplecticPacking.Bundle.cF_cube, SymplecticPacking.Bundle.dF_sq,
SymplecticPacking.Bundle.torusLift_fermat, SymplecticPacking.Bundle.torusPoint_mem,
SymplecticPacking.Bundle.wp_injective, SymplecticPacking.Bundle.torusLift_proj_injective,
SymplecticPacking.Bundle.torusPoint_proj_injective}
\uses{def:fer-torus, lem:fer-curve}
The invariants of $\Lambda_0$ are $g_2=0$ and $g_3\ne0$, so $w_\wp^2=4\wp^3-g_3$. The map $\Phi_\Gamma$ takes values in gauge classes of nonzero Fermat vectors, and it is injective on $\C/\Lambda_0$.
\end{lemma}

\begin{proof}
\emph{Invariants.} The Eisenstein series satisfy $G_k(\lambda\Lambda)=\lambda^{-k}G_k(\Lambda)$. With $\lambda=\zeta_6$ and
$\zeta_6\Lambda_0=\Lambda_0$ this gives $G_4=\zeta_6^{-4}G_4$, and $\zeta_6^4\ne1$, so $G_4=0$ and $g_2=60G_4=0$. Suppose
$g_3=0$, so $w_\wp^2=4\wp^3$. At the half period $\omega_1=\frac12$, $\wp'(\omega_1)=0$ ($\wp'$ is odd and periodic), hence
$\wp(\omega_1)=0$. If $\wp-\wp(\omega_1)$ vanishes to order $n\ge1$ at $\omega_1$, then $\wp'$ vanishes to order $n-1$, and
$w_\wp^2=4\wp^3$ gives $2(n-1)=3n$, which is impossible. So $g_3\ne0$, and $c_F,d_F$ are defined and nonzero.

\emph{Image.} Off $\Lambda_0$,
\[
 (2c_F\wp)^3-(1+d_Fw_\wp)^3-(1-d_Fw_\wp)^3=8c_F^3\wp^3-2-6d_F^2w_\wp^2=\frac{8\wp^3-2g_3-2w_\wp^2}{g_3}=0 .
\]
At the lattice point, $0^3+1^3+(-1)^3=0$. At the lattice point the formula takes this value.

\emph{Injectivity.} For $z\notin\Lambda_0$ the last two homogeneous coordinates of $\Phi_\Gamma(z)$ sum to $-2\ne0$, while
for $[0:1:-1]$ they sum to $0$; so the lattice point is the only preimage of $[0:1:-1]$. If $\Phi_\Gamma(z)=\Phi_\Gamma(z')$ with
$z,z'\notin\Lambda_0$, the representatives agree up to a scalar, which is $1$ by the sums of the last two coordinates;
since $c_F,d_F\ne0$, $\wp(z)=\wp(z')$ and $\wp'(z)=\wp'(z')$. The elliptic function $\wp-\wp(z)$ has a single pole in
$\C/\Lambda_0$, of order two, so it has exactly two zeros counted with multiplicity; since $\wp$ is even, these are $z$ and
$-z$. Hence $z'\equiv\pm z$ modulo $\Lambda_0$. If $z'\equiv-z$, then $\wp'(z)=\wp'(z')=-\wp'(z)$, so $\wp'(z)=0$. The odd
elliptic function $\wp'$ has a single pole, of order three, so it has exactly three zeros counted with multiplicity, and it
vanishes at the three half periods; so $z$ is a half period and $-z\equiv z$. In either case $z'\equiv z$.
\end{proof}

\begin{definition}[Theta lift]\label{def:fer-theta}
\lean{SymplecticPacking.Bundle.jacobiTheta1, SymplecticPacking.Bundle.thetaFactor,
SymplecticPacking.Bundle.rTheta, SymplecticPacking.Bundle.wpConst,
SymplecticPacking.Bundle.PTheta, SymplecticPacking.Bundle.QTheta,
SymplecticPacking.Bundle.Zraw, SymplecticPacking.Bundle.Zunit}
\uses{def:fer-torus}
Let $x=\pi z=\pi(u+\zeta v)$ and
\[
 \vartheta(x)=2\sum_{n\ge0}(-1)^ne^{i\pi\zeta(n+1/2)^2}\sin\big((2n+1)x\big),
\]
the Jacobi function $\vartheta_1(x\,|\,\zeta)$; primes denote $x$-derivatives. Put
$r_\vartheta=\vartheta'''(0)/(3\vartheta'(0))$,
\[
 P_\vartheta=\vartheta'^2\vartheta-\vartheta''\vartheta^2+r_\vartheta\vartheta^3,\qquad
 Q_\vartheta=-\vartheta'''\vartheta^2+3\vartheta'\vartheta''\vartheta-2\vartheta'^3,
\]
\[
 \hat Z=\big(2c_F\pi^2P_\vartheta,\;-\vartheta^3-d_F\pi^3Q_\vartheta,\;-\vartheta^3+d_F\pi^3Q_\vartheta\big),\qquad
 Z=\frac{\hat Z}{\|\hat Z\|},
\]
all evaluated at $x=\pi z$.
\end{definition}

Equivalently, $\vartheta(x)=-i\,e^{i\pi\zeta/4+ix}\,\Theta(x/\pi+\frac12+\frac\zeta2,\zeta)$ with
$\Theta(y,\zeta)=\sum_{n\in\Z}e^{2\pi iny+i\pi n^2\zeta}$.

\begin{lemma}[The lift has no zero]\label{lem:fer-theta}
\lean{SymplecticPacking.Bundle.jacobiTheta1_pi_mul_eq_zero_iff,
SymplecticPacking.Bundle.dTheta1_ne_zero_of_eq_zero, SymplecticPacking.Bundle.wp_eq_theta,
SymplecticPacking.Bundle.Zraw_eq, SymplecticPacking.Bundle.Zraw_ne_zero,
SymplecticPacking.Bundle.Zraw_automorphy, SymplecticPacking.Bundle.contDiff_Zunit}
\uses{def:fer-theta, lem:fer-torus}
The zeros of $\vartheta$ are exactly the points of $\pi\Lambda_0$, and they are simple. Off $\Lambda_0$,
\[
 \wp(z)=-\pi^2(\log\vartheta)''(\pi z)+\pi^2r_\vartheta ,
\]
and $\hat Z=\vartheta^3\cdot(2c_F\wp,\,-1-d_Fw_\wp,\,-1+d_Fw_\wp)$ wherever $\vartheta\ne0$. $\hat Z$ is an entire function
of $z$ without zeros, so $Z$ is a smooth map $\R^2\to\C^3$ with $\|Z\|=1$. Under $(u,v)\mapsto(u+1,v)$ and
$(u,v)\mapsto(u,v+1)$, $\hat Z$ is multiplied by $-1$ and by $\mathfrak m(\pi z)^3$, where
$\mathfrak m(x)=-e^{-i\pi\zeta}e^{-2ix}$.
\end{lemma}

\begin{proof}
\emph{Quasi-periodicity.} From the series, $\vartheta$ is entire and odd, so $\vartheta(0)=0$, and
$\vartheta(x+\pi)=-\vartheta(x)$, $\vartheta(x+\pi\zeta)=\mathfrak m(x)\vartheta(x)$ (these follow from the corresponding
identities of $\Theta$). Since $\mathfrak m'/\mathfrak m=-2i$, the logarithmic derivative $\ell=\vartheta'/\vartheta$ satisfies
$\ell(x+\pi)=\ell(x)$ and $\ell(x+\pi\zeta)=\ell(x)-2i$.

\emph{Zeros.} For every $q\in\C$ there is a parallelogram $P=a+[0,1]\pi+[0,1]\pi\zeta$ containing $q$ in its interior with no
zero of $\vartheta$ on its boundary. By the argument principle the number of zeros in $P$ is
$(2\pi i)^{-1}\oint_{\partial P}\ell\,dx$. The two edges parallel to $\pi\zeta$ cancel, and the two edges parallel to $\pi$
contribute $\int_a^{a+\pi}(\ell(x)-\ell(x+\pi\zeta))\,dx=2\pi i$. So every such $P$ contains exactly one zero, and it is
simple. Since $\vartheta(0)=0$ and the zero set is $\pi\Lambda_0$-invariant, every point of $\pi\Lambda_0$ is a zero. Conversely, let
$q$ be a zero and $P\ni q$ as above. Then $\partial P$ contains no point of $\pi\Lambda_0$, so the translate $P$ of a fundamental
parallelogram contains exactly one point of $\pi\Lambda_0$, in its interior; this point and $q$ are zeros in $P$, so they coincide.
Thus the zeros are exactly $\pi\Lambda_0$; in particular $\vartheta'\ne0$ there.

\emph{The formula for $\wp$.} Put $F(z)=-\pi^2(\log\vartheta)''(\pi z)=\pi^2(\vartheta'^2-\vartheta\vartheta'')/\vartheta^2$. Since $\ell$
changes only by additive constants under $\pi\Lambda_0$, $F$ is $\Lambda_0$-periodic. From
$\vartheta(x)=\vartheta'(0)x\big(1+\frac{\vartheta'''(0)}{6\vartheta'(0)}x^2+O(x^4)\big)$ we get
$(\log\vartheta)''=-x^{-2}+\frac{\vartheta'''(0)}{3\vartheta'(0)}+O(x^2)$, so $F(z)=z^{-2}-\pi^2r_\vartheta+O(z^2)$. Thus $\wp-F$ is
periodic, holomorphic off $\Lambda_0$ and bounded near the lattice, hence constant by Liouville's theorem; the constant is
$\pi^2r_\vartheta$ because $\wp$ has zero constant Laurent term at $0$. So $\vartheta^3\wp=\pi^2P_\vartheta$. Differentiating in $z$
(one more factor $\pi$),
$w_\wp=\pi^3\frac{d}{dx}\frac{\vartheta'^2-\vartheta\vartheta''}{\vartheta^2}=\pi^3Q_\vartheta/\vartheta^3$, which gives the formula for
$\hat Z$.

\emph{No zero.} Where $\vartheta\ne0$ the sum of the last two components is $-2\vartheta^3\ne0$. Where $\vartheta=0$,
$\vartheta'\ne0$ and $Q_\vartheta=-2\vartheta'^3$, so the second component is $2d_F\pi^3\vartheta'^3\ne0$. $\hat Z$ is entire,
being a polynomial in $\vartheta,\dots,\vartheta'''$ at $\pi z$.

\emph{Automorphy.} Under $x\mapsto x+\pi$ all of $\vartheta,\dots,\vartheta'''$ change sign, and $P_\vartheta$, $Q_\vartheta$,
$\vartheta^3$ are homogeneous of degree three in them, so $\hat Z\mapsto-\hat Z$. Under $x\mapsto x+\pi\zeta$, off
$\pi\Lambda_0$ we have $P_\vartheta=\vartheta^3\wp/\pi^2$ and $Q_\vartheta=\vartheta^3w_\wp/\pi^3$ by the previous step; since $\wp$ and
$w_\wp$ are $\Lambda_0$-periodic, $P_\vartheta$, $Q_\vartheta$ and $\vartheta^3$ are all multiplied by $\mathfrak m(x)^3$. By continuity
the same holds at the points of $\pi\Lambda_0$.
\end{proof}

At the lattice points $P_\vartheta=0$ and $Q_\vartheta=-2\vartheta'^3$, so $[\hat Z]=[0:1:-1]$; off the lattice
$[\hat Z]=\Phi_\Gamma(z)$. Hence $Z$ is a smooth unit Fermat lift on $\R^2$ with that gauge class.

\begin{construction}[The atlas of $N$]\label{cons:fer-atlas}
\lean{SymplecticPacking.Return.Ncoord, SymplecticPacking.Embed.deckU, SymplecticPacking.Embed.deckV,
SymplecticPacking.Return.betaLoc, SymplecticPacking.Return.d1_betaLoc, SymplecticPacking.Bundle.Omega0loc}
\uses{lem:fer-theta, def:cub-model}
\begin{enumerate}[label=(\alph*)]
\item \emph{Base.} The chart $(u,v)\mapsto\Phi_\Gamma(u+\zeta v)$ is the injective Fermat representative map of the base torus $\Gamma=(\R/\Z)^2$; all base objects are written in $(u,v)$, and $\nu=du\wedge dv$.
\item \emph{The torus cover of $N$.} The map $\R^2\times\C\to N$, $(u,v,w)\mapsto[Z(u,v),w]$, with $w=\xi+i\eta$, is the cover-to-representative map. Let $\phi(u,v)=\mathfrak m(\pi z)^3/|\mathfrak m(\pi z)|^3$, so that $Z(u+1,v)=-Z(u,v)$
and $Z(u,v+1)=\phi(u,v)Z(u,v)$ (Lemma~\ref{lem:fer-theta}). The deck maps
\[
 T_u(u,v,w)=(u+1,v,-w),\qquad T_v(u,v,w)=(u,v+1,\phi(u,v)^3w)
\]
generate a free $\Z^2$-action, and $(u,v,w)\mapsto[Z(u,v),w]$ is exactly the deck relation defining $N$.
\item \emph{Algebraic frames.} For $j=0,1,2$, $Z\bar Z_j/|Z_j|$ is a local unitary frame over $\Gamma\cap\{Z_j\ne0\}$ in the sense
of Definition~\ref{def:cub-model}. Its fibre coordinate $w_j$ is given by $U=(Z_j/|Z_j|)^3w_j$, so its unit section is
$[Z]\mapsto[Z,(Z_j/|Z_j|)^3]$, and its connection form is $-3\eta_H+\frac3{2\pi}d\arg Z_j$. These three sets cover $\Gamma$, and the
transition functions $w_j=(Z_k|Z_j|/(|Z_k|Z_j))^3w_k$ have weight three.
\item \emph{The model form in the torus frame.} In the frame (b), $p=\pi|w|^2$, $\alpha=d\theta-3Z^*\eta_H$, and
\[
 \Omega_0=d\Big((1-3p)\,Z^*\eta_H+\tfrac12(\xi\,d\eta-\eta\,d\xi)\Big)=\omega_V+(1-3p)\sigma+dp\wedge(-3Z^*\eta_H).
\]
\end{enumerate}
\end{construction}

\begin{proof}[Verification]
(b) $[Z(u+1,v),-w]=[-Z(u,v),-w]=[Z(u,v),(-1)^{-3}(-w)]=[Z(u,v),w]$ by $(e^{i\pi}Z,e^{3i\pi}U)\sim(Z,U)$, and similarly
$[Z(u,v+1),\phi^3w]=[\phi Z(u,v),\phi^3w]=[Z(u,v),w]$; so the map is invariant. Since $\mathfrak m(x+\pi)=\mathfrak m(x)$, we have
$\phi(u+1,v)=\phi(u,v)$, so $T_uT_v(u,v,w)=(u+1,v+1,-\phi(u,v)^3w)=T_vT_u(u,v,w)$. If two unit representative pairs from the cover are gauge equivalent, projective injectivity of the explicit
torus parametrization makes their base coordinates differ by an integral lattice vector.
Apply the corresponding deck map; the base coordinates then agree and unit gauge equivalence fixes the fibre coordinate.
Thus gauge equivalence on this cover is exactly the stated deck relation.
(c) $Z\bar Z_j/|Z_j|$ is unchanged under $Z\mapsto e^{i\psi}Z$, so it is a gauge-invariant normalized lift on the cover patch $\{Z_j\ne0\}$, and
$[Z,(Z_j/|Z_j|)^3w_j]=[e^{-i\arg Z_j}Z,w_j]$ by the weight-three relation. Since $\arg U=3\arg Z_j+\arg w_j$, the angular
form of Lemma~\ref{lem:fer-alpha} is $\alpha=d\arg w_j/(2\pi)+\frac3{2\pi}d\arg Z_j-3\eta_H$, so the connection form of the
frame is $-3\eta_H+\frac3{2\pi}d\arg Z_j$. No point of $\Gamma$ has
$Z_0=Z_1=Z_2=0$. On overlaps, $[Z,(Z_k/|Z_k|)^3w_k]=[Z,(Z_j/|Z_j|)^3w_j]$ gives $w_j=(Z_k|Z_j|/(|Z_k|Z_j))^3w_k$.
(d) Lemma~\ref{lem:fer-alpha} and Lemma~\ref{lem:fer-Omega0} with $a_{\rm loc}=-3Z^*\eta_H$, and $d(Z^*\eta_H)=\sigma$; the
primitive is $Z^*\eta_H+p\alpha$.
\end{proof}

\begin{lemma}[The deck phase]\label{lem:fer-clutch}
\lean{SymplecticPacking.Embed.deckU,SymplecticPacking.Embed.deckV,SymplecticPacking.Bundle.Zraw_automorphy}
The two fibre factors are $-1$ and $\phi^3$.
They commute with the lattice translations. On $v=0$,
$\phi(u,0)=\phi_0e^{-6\pi iu}$ for a constant $|\phi_0|=1$.
\end{lemma}
\begin{proof}
The theta multiplier is periodic under $u\mapsto u+1$, so the two deck maps commute.
On $v=0$, its formula is $\mathfrak m(\pi u)=-e^{-i\pi\zeta}e^{-2\pi iu}$.
Taking its unit phase and cubing gives the displayed identity.
\end{proof}

\begin{lemma}[The area of the base]\label{lem:fer-sigma}
\lean{SymplecticPacking.Bundle.d_etaH, SymplecticPacking.Bundle.sigmaC_pos,
SymplecticPacking.Bundle.sigmaC_periodic, SymplecticPacking.Bundle.integral_sigma}
\uses{lem:fer-theta}
Write $Z^*\eta_H=\mathsf h_u\,du+\mathsf h_v\,dv$ and $\sigma=d(Z^*\eta_H)=\sigma_\Gamma(u,v)\,du\wedge dv$ on $\R^2$. Then
$\sigma_\Gamma$ is smooth, $1$-periodic in $u$ and in $v$, strictly positive, and
$\int_0^1\!\int_0^1\sigma_\Gamma\,du\,dv=3$.
\end{lemma}

\begin{proof}
\emph{Positivity.} $\sigma_\Gamma$ is the pull-back of $\omega_{\rm FS}=dd^c\log\|Z\|^2$ by $z\mapsto[\hat Z(z)]$:
\[
 \sigma_\Gamma=\operatorname{Im}\zeta\cdot\frac{\|\hat Z\|^2\|\hat Z'\|^2-|\hat Z^\dagger \hat Z'|^2}{\pi\|\hat Z\|^4}
 \qquad(\hat Z'=d\hat Z/dz),
\]
which by the Cauchy--Schwarz equality case is positive unless $\hat Z'$ is proportional to $\hat Z$. Off the lattice,
$\hat Z=\vartheta^3R$ with $R=(2c_F\wp,-1-d_Fw_\wp,-1+d_Fw_\wp)$, and $\hat Z'\parallel \hat Z$ iff $R'\parallel R$. Now
$R'=(2c_Fw_\wp,-d_F\wp'',d_F\wp'')$; if $R'=aR$, adding the last two components gives $0=-2a$, so $R'=0$, i.e.\ $w_\wp=0$ and
$\wp''=0$. This is impossible: differentiating $w_\wp^2=4\wp^3-g_3$ gives $\wp''=6\wp^2$, and $w_\wp=0$ forces $4\wp^3=g_3\ne0$.
At a lattice point the first component of $\hat Z$ vanishes while that of $\hat Z'$ is
$2c_F\pi^3P_\vartheta'=2c_F\pi^3\vartheta'^3\ne0$, so again $\hat Z'\not\parallel \hat Z\ne0$.

\emph{Periodicity.} By Lemma~\ref{lem:fer-theta}, $Z$ changes by unit phases $\phi$ under the deck translations, and
$\eta_H$ changes by the closed form $(2\pi i)^{-1}\bar\phi\,d\phi$; so $\sigma=d(Z^*\eta_H)$ is periodic.

\emph{Integral.} By Stokes on the unit square, $\iint\sigma_\Gamma=\iint(\partial_u\mathsf h_v-\partial_v\mathsf h_u)$. Under $u\mapsto u+1$
the lift changes by the constant phase $-1$, so $\mathsf h_v(1,v)=\mathsf h_v(0,v)$ and the first term integrates to $0$. Under
$v\mapsto v+1$ it changes by $\phi(u,v)$, and on the edge $v=0$, $\phi=\phi_0e^{-6\pi iu}$ (Lemma~\ref{lem:fer-clutch}); hence
$\mathsf h_u(u,1)=\mathsf h_u(u,0)+(2\pi i)^{-1}(-6\pi i)=\mathsf h_u(u,0)-3$. So
$\iint\partial_v\mathsf h_u=\int_0^1(\mathsf h_u(u,1)-\mathsf h_u(u,0))\,du=-3$ and $\iint\sigma_\Gamma=3$.
\end{proof}

\begin{lemma}[The three points on $L$]\label{lem:fer-bpoints}
\lean{SymplecticPacking.Bundle.bpoints, SymplecticPacking.Bundle.Z0_eq_zero_iff}
\uses{lem:fer-theta}
The zero set $\mathcal Z=\{(u,v)\in\Gamma:Z_0(u,v)=0\}$, the set of base points lying on the line $L=\{Z_0=0\}$,
consists of the three points
\[
 b_1=(0,0),\qquad b_2=(\tfrac13,\tfrac23),\qquad b_3=(\tfrac23,\tfrac13).
\]
\end{lemma}

\begin{proof}
$Z_0=0$ iff $P_\vartheta=0$, i.e.\ iff $\vartheta=0$ (the point $b_1$, where $\Phi_\Gamma=[0:1:-1]$) or $\wp=0$. Since $\zeta\Lambda_0=\Lambda_0$,
$\wp(\zeta z)=\zeta^{-2}\wp(z)=\zeta\wp(z)$; at a non-lattice fixed point of $z\mapsto\zeta z$ modulo $\Lambda_0$ this forces
$\wp(z)=0$. For $z_2=\frac13+\frac23\zeta$, $\zeta z_2=-\frac23-\frac13\zeta\equiv z_2$ (using $\zeta^2=-1-\zeta$), and similarly for
$z_3=\frac23+\frac13\zeta\equiv-z_2$. The elliptic function $\wp$ has a single pole, of order two, in $\C/\Lambda_0$,
hence exactly two zeros counted with multiplicity; since $z_2\not\equiv z_3$, these are $z_2$ and $z_3$.
\end{proof}

\subsection{Handle data}\label{sec:fer-handle}

The handle data of the cubic model are stated in Proposition~\ref{prop:cub-handle}. Picture the torus as the unit
square in $(u,v)$. The annulus $\mathcal A$ is the vertical band $\frac18\le u\le\frac5{24}$ inside the wider band
$\mathcal A^+=\{\frac1{12}\le u\le\frac14\}$; the dual band $\mathcal B=\{\frac19<v<\frac29\}$ contains the closed clock band
$\frac18\le v\le\frac5{24}$ and carries the closed one-form $\beta=b(v)\,dv$; the two bands meet in one square. The three points $b_j$ lie far from both
bands and carry small discs $U_j$. This subsection gives the frame over $\mathcal A^+$, the primitive operator and the
cutoffs used in \S\ref{sec:fer-alpha1}.

\begin{definition}[Frame over $\mathcal A^+$]\label{def:fer-frame}
\lean{SymplecticPacking.Bundle.frameA, SymplecticPacking.Handle.chartInv,
SymplecticPacking.Bundle.lambdaC, SymplecticPacking.Bundle.aA}
\uses{cons:fer-atlas}
On $\Gamma\cap\{Z_0\ne0\}$ let $w_{\mathcal A}$ be the fibre coordinate of the local unitary frame $Z\bar Z_0/|Z_0|$ of
Construction~\ref{cons:fer-atlas}(c), i.e.\ $U=(Z_0/|Z_0|)^3w_{\mathcal A}$; we also call it the frame $(Z_0/|Z_0|)^3$, after its unit
section. Let $\theta_{\mathcal A}=\arg w_{\mathcal A}/2\pi$, and put
\[
 \lambda_\Gamma=\eta_H-\frac{1}{2\pi}\operatorname{Im}\frac{dZ_0}{Z_0},\qquad a_{\mathcal A}=-3\lambda_\Gamma,\qquad
 \Gamma_{\mathcal A}(p)=(1-3p)\,\lambda_\Gamma ,
\]
so that $a_{\mathcal A}$ is the connection form $a_{\rm loc}$ of this frame.
\end{definition}

\begin{lemma}[The frame over the annulus]\label{lem:fer-frame}
\lean{SymplecticPacking.Bundle.lambdaC_descends, SymplecticPacking.Bundle.d_lambdaC,
SymplecticPacking.Bundle.angularForm_frameA}
\uses{def:fer-frame, lem:fer-alpha}
$\lambda_\Gamma$ is a smooth one-form on $\Gamma\cap\{Z_0\ne0\}$ (no branch of $\arg$ is needed), $d\lambda_\Gamma=\sigma$, and on
$N|_{\Gamma\cap\{Z_0\ne0\}}$
\[
 \alpha=d\theta_{\mathcal A}+a_{\mathcal A},\qquad \Omega_0=\omega_V+d\Gamma_{\mathcal A},\quad \omega_V=dp\wedge d\theta_{\mathcal A}.
\]
\end{lemma}

\begin{proof}
$\operatorname{Im}(dZ_0/Z_0)=d\arg Z_0$ is a smooth closed form where $Z_0\ne0$. The form $\lambda_\Gamma$ on $S_\Gamma$ is invariant and
vanishes on the generator $iZ$ ($\frac1{2\pi}-\frac1{2\pi}$), so it is basic; $d\lambda_\Gamma=d\eta_H=\sigma$. Since
$\arg U=3\arg Z_0+\arg w_{\mathcal A}$,
$\alpha=d\theta_{\mathcal A}+\frac{3}{2\pi}d\arg Z_0-3\eta_H=d\theta_{\mathcal A}-3\lambda_\Gamma$. Then
$\Omega_0=d\lambda_\Gamma+d(p\,d\theta_{\mathcal A})-3\,d(p\lambda_\Gamma)=\omega_V+d\big((1-3p)\lambda_\Gamma\big)$.
\end{proof}

\begin{definition}[Torus primitive operator]\label{def:fer-Pprim}
\lean{SymplecticPacking.Bundle.Pprim, SymplecticPacking.Bundle.Pb0,
SymplecticPacking.Bundle.PA0, SymplecticPacking.Bundle.PB0}
This is the operator of Proposition~\ref{prop:cub-handle}(e). For a smooth function $\mathfrak b$ on $(\R/\Z)^2$ with
$\int\mathfrak b\,\nu=0$ put
\[
 \mathfrak b_0(v)=\int_0^1\mathfrak b(s',v)\,ds',\quad
 A_0(u,v)=\int_0^u\big(\mathfrak b(s',v)-\mathfrak b_0(v)\big)\,ds',\quad
 B_0(v)=\int_0^v\mathfrak b_0(t)\,dt,
\]
and $P_\Gamma(\mathfrak b\,\nu)=A_0\,dv-B_0\,du$.
\end{definition}

\begin{lemma}\label{lem:fer-Pprim}
\lean{SymplecticPacking.Bundle.Pprim_periodic, SymplecticPacking.Bundle.d_Pprim,
SymplecticPacking.Bundle.contDiff_Pprim_cu, SymplecticPacking.Bundle.contDiff_Pprim_cv,
SymplecticPacking.Bundle.Pprim_affine}
\uses{def:fer-Pprim}
$A_0$ and $B_0$ are $1$-periodic in both variables, so $P_\Gamma(\mathfrak b\nu)$ is a smooth one-form on the torus, and
$d_\Sigma P_\Gamma(\mathfrak b\nu)=\mathfrak b\nu$. $P_\Gamma$ is linear, and if $\mathfrak b$ depends smoothly on parameters then so does
$P_\Gamma(\mathfrak b\nu)$, with all jets; in particular it preserves affine dependence on $p$.
\end{lemma}

\begin{proof}
$A_0(u+1,v)-A_0(u,v)=\int_u^{u+1}(\mathfrak b-\mathfrak b_0)\,ds'=0$, and $A_0$ is periodic in $v$ because $\mathfrak b$ and $\mathfrak b_0$
are; $B_0(v+1)-B_0(v)=\int_0^1\mathfrak b_0=\int\mathfrak b\,\nu=0$. Then
$d(A_0\,dv-B_0\,du)=\partial_uA_0\,du\wedge dv-\mathfrak b_0\,dv\wedge du=(\mathfrak b-\mathfrak b_0+\mathfrak b_0)\,\nu$. Smooth dependence
on parameters follows by differentiation under the integral sign.
\end{proof}

\begin{proposition}[Cutoffs on the cubic handle]\label{prop:fer-handle}
\lean{SymplecticPacking.Bundle.annulusPlus, SymplecticPacking.Bundle.annulus,
SymplecticPacking.Bundle.psiB, SymplecticPacking.Bundle.bB, SymplecticPacking.Bundle.betaB,
SymplecticPacking.Bundle.stripB, SymplecticPacking.Bundle.zetaA, SymplecticPacking.Bundle.zetaJ,
SymplecticPacking.Bundle.supports, SymplecticPacking.Bundle.Zunit_zero_ne_zero_of_mem}
\uses{prop:cub-handle, lem:fer-bpoints, def:pre-cutoff, lem:pre-cutoff}
Let $\mathcal A\subset\mathcal A^+$, $\mathcal B$ and $\beta$ be the handle data of Proposition~\ref{prop:cub-handle}, and take
the discs of its item (g) closed: $U_j=\{x_j^2+y_j^2\le10^{-4}\}$, of radius $\frac1{100}$ about $b_j$ in the centred coordinates
$(x_j,y_j)=(u-u(b_j),v-v(b_j))$, reduced to $[-\frac12,\frac12)^2$. With the flat step $S$ and the cutoffs of
Definition~\ref{def:pre-cutoff}, the cutoffs
\[
 \zeta_{\mathcal A}(u)=\chi^{\rm I}_{9/96,\,23/96,\,2/96}(u),\qquad \zeta_j=\chi_{1/200,\,1/100}(x_j,y_j)
\]
(extended by zero) satisfy: $\zeta_{\mathcal A}=1$ on $\{\frac{11}{96}\le u\le\frac{21}{96}\}$, a neighbourhood of $\mathcal A$, and
$\operatorname{supp}\zeta_{\mathcal A}=\{\frac9{96}\le u\le\frac{23}{96}\}\subset\operatorname{int}\mathcal A^+$; $\zeta_j=1$ on the disc
of radius $\frac1{200}$ about $b_j$, and $\operatorname{supp}\zeta_j=U_j$. The discs $U_j$ are pairwise disjoint and disjoint from
$\mathcal A^+\cup\mathcal B$, and $\operatorname{supp}\beta\cap U_j=\varnothing$. All these data are independent of the input.
\end{proposition}

\begin{proof}
By Lemma~\ref{lem:pre-cutoff}(b), $\chi^{\rm I}_{9/96,23/96,2/96}=1$ on $[\frac{11}{96},\frac{21}{96}]\supset
[\frac{12}{96},\frac{20}{96}]=\mathcal A$, and it is positive exactly on $(\frac9{96},\frac{23}{96})$, so its support is
$[\frac9{96},\frac{23}{96}]$; moreover $\frac8{96}<\frac9{96}$ and $\frac{23}{96}<\frac{24}{96}$, so this interval lies in the interior
of $\mathcal A^+=[\frac8{96},\frac{24}{96}]$. By Lemma~\ref{lem:pre-cutoff}(a), $\chi_{1/200,1/100}=1$ on the closed disc of radius
$\frac1{200}$, vanishes outside the open disc of radius $\frac1{100}$, and is positive on that open disc, since $S<1$ on $(-\infty,1)$;
so its support is the closed disc $U_j$. The $b_j$ have mutual distance at least $\frac13$ in the flat metric of the torus,
and distance at least $\frac1{12}$ from $\mathcal A^+\cup\mathcal B$ (proof of Proposition~\ref{prop:cub-handle}(a),(b)); both
bounds exceed $\frac2{100}$, which gives the disjointness. Finally $\operatorname{supp}\beta\subset\mathcal B$.
\end{proof}

%----------------------------------------------------------------------
\subsection{A global primitive vanishing on three whole fibres}\label{sec:fer-alpha1}

We now write the first path of \Mref{Proposition}{prop:hdl-first-path} for the cubic in closed form. We start from an explicit
primitive $\alpha^0$ of the exact base form $\kappa_1-\kappa$ and correct it twice by $d_\Sigma$-closed forms, which leaves the base
curvature unchanged: over $\mathcal A^+$ we remove the difference between $\alpha^0$ and the desired local primitive (a closed
form on an annulus, exact after subtracting its period), and near the three points $b_j$ we remove the value of the
primitive. The supports of the two corrections are disjoint (Proposition~\ref{prop:fer-handle}).

The scalar data in this subsection are the concrete formulas below. They apply to both branches of the four-dimensional comparator.

\begin{definition}[Admissible auxiliary data and base area forms]\label{def:fer-kappa}
\lean{SymplecticPacking.Branch2.CubicData.eh,
SymplecticPacking.Branch2.CubicData.Hstar,
SymplecticPacking.Bundle.kappa,
SymplecticPacking.Bundle.Cubic.kappa1,
SymplecticPacking.Bundle.chiA}
\uses{prop:fer-handle, M-eq:hdl-eh}
Let $0<h<1/3$ and $0<g\le3(1-3h)$. Set
\[
 e_h=g/16,\qquad H^*(p)=3(1-3p)-e_h,\qquad
 f(s)=S(2s-1/2),\qquad \Gamma_{\rm sm}=e_h(u-1/6)\,dv.
\]
Thus $0<e_h\le\Hb(h)/2$, $f$ is nondecreasing, zero for $s\le1/4$ and one for $s\ge3/4$,
and $d\Gamma_{\rm sm}=e_h\nu$.
Put $\chi=f'(s)\,ds\wedge\beta$ (on $\mathcal A$, extended by zero) and
\[
 \kappa(p)=(1-3p)\,\sigma,\qquad \kappa_1(p)=e_h\nu+H^*(p)\chi .
\]
\end{definition}

\begin{lemma}[Base area forms]\label{lem:fer-kappa}
\lean{SymplecticPacking.Bundle.integral_chi,
SymplecticPacking.Bundle.integral_kappa,
SymplecticPacking.Bundle.Cubic.integral_kappa1,
SymplecticPacking.Bundle.Cubic.kappa_pos}
\uses{def:fer-kappa, lem:fer-sigma, prop:fer-handle}
$\chi\ge0$, $\int\chi=1$, and $\int\kappa(p)=\int\kappa_1(p)=\Hb(p)$. On $[0,h]$, $H^*\ge\frac12\Hb(h)>0$. For every $p<\frac13$
with $H^*(p)\ge0$, both $\kappa(p)$ and $\kappa_1(p)$ are positive area forms, and $\kappa_1(p)\ge e_h\nu$.
\end{lemma}

\begin{proof}
$\chi=f'(s)b(v)\,ds\wedge dv\ge0$ since $f'\ge0$, $b\ge0$, and $\int\chi=\int_0^1f'\,ds\cdot\int_0^1b\,dv=1$. By
Lemma~\ref{lem:fer-sigma}, $\int\kappa=3(1-3p)$, and $\int\kappa_1=e_h+H^*=\Hb$. On $[0,h]$, $\Hb\ge\Hb(h)$ and
$e_h\le\frac12\Hb(h)$. Finally $\sigma>0$ (Lemma~\ref{lem:fer-sigma}), and $\kappa_1-e_h\nu=H^*\chi\ge0$.
\end{proof}

\begin{definition}[The annular difference and its period]\label{def:fer-Delta}
\lean{SymplecticPacking.Bundle.Cubic.alpha10,
SymplecticPacking.Bundle.Cubic.gammaSm,
SymplecticPacking.Bundle.Cubic.DeltaA,
SymplecticPacking.Bundle.Cubic.corePeriod}
\uses{def:fer-kappa, def:fer-Pprim, def:fer-frame}
Let $\alpha^0(p)=P_\Gamma\big(\kappa_1(p)-\kappa(p)\big)$, a base one-form with coefficients depending on the base point and
affinely on $p$, with $d_\Sigma\alpha^0=\kappa_1-\kappa$ (Lemmas~\ref{lem:fer-Pprim} and~\ref{lem:fer-kappa}). On $\mathcal A^+$ put
\[
 \Delta_0(p)=\Gamma_{\mathcal A}(p)+\alpha^0(p)-\Gamma_{\rm sm}-f(s)H^*(p)\beta ,
\]
and let $c(p)=\oint_{\{u={\rm const}\}}\Delta_0(p)$ be its period over a $v$-circle in $\mathcal A^+$.
\end{definition}

\begin{lemma}\label{lem:fer-Delta}
\lean{SymplecticPacking.Bundle.Cubic.d_alpha10,
SymplecticPacking.Bundle.Cubic.alpha10_affine,
SymplecticPacking.Bundle.Cubic.alpha10_periodic,
SymplecticPacking.Bundle.Cubic.DeltaA_closed,
SymplecticPacking.Bundle.Cubic.corePeriod_indep,
SymplecticPacking.Bundle.Cubic.corePeriod_affine}
\uses{def:fer-Delta, lem:fer-frame}
For each $p$, $\Delta_0(p)$ is $d_\Sigma$-closed on $\mathcal A^+$; hence $c(p)$ does not depend on the circle, and $c$ is affine
in $p$.
\end{lemma}

\begin{proof}
$d_\Sigma\Gamma_{\mathcal A}=(1-3p)\sigma=\kappa$ by Lemma~\ref{lem:fer-frame}, $d_\Sigma\alpha^0=\kappa_1-\kappa$,
$d_\Sigma\Gamma_{\rm sm}=e_h\nu$ and $d_\Sigma(fH^*\beta)=H^*f'\,ds\wedge\beta=H^*\chi$; so
$d_\Sigma\Delta_0=\kappa+\kappa_1-\kappa-e_h\nu-H^*\chi=0$. Two $v$-circles bound an annulus in $\mathcal A^+$, so by Stokes the period
does not depend on the circle. Every term of $\Delta_0$ is affine in $p$ ($\Hb$ is affine, $e_h$ and $\Gamma_{\rm sm}$ are constant in
$p$, and $P_\Gamma$ is linear).
\end{proof}

The period $c(p)$ records the holonomy of the frame $(Z_0/|Z_0|)^3$ around the core circle together with the periods of the
chosen primitives; in general $c(p)\ne0$. Since $\oint\beta=1$, $\Delta_0-c(p)\beta$ has period zero and is $d_\Sigma$-exact on
$\mathcal A^+$.

\begin{definition}[Annular primitive]\label{def:fer-FA}
\lean{SymplecticPacking.Bundle.Cubic.FA}
\uses{lem:fer-Delta}
Write $\Delta_0(p)-c(p)\beta=D_u\,du+D_v\,dv$ on $\mathcal A^+$ and put
\[
 F_{\mathcal A}(u,v,p)=\int_{1/6}^uD_u(t,0,p)\,dt+\int_0^vD_v(u,t,p)\,dt\qquad(\tfrac1{12}\le u\le\tfrac14,\ v\in\R),
\]
the integral of $\Delta_0-c\beta$ along the path from $(\frac16,0)$ horizontally and then vertically.
\end{definition}

\begin{lemma}\label{lem:fer-FA}
\lean{SymplecticPacking.Bundle.Cubic.FA_periodic,
SymplecticPacking.Bundle.Cubic.d_FA,
SymplecticPacking.Bundle.Cubic.contDiffOn_FA,
SymplecticPacking.Bundle.Cubic.FA_affine}
\uses{def:fer-FA, prop:cub-handle}
$F_{\mathcal A}$ is a smooth function on $\mathcal A^+$, $1$-periodic in $v$, affine in $p$, and
$d_\Sigma F_{\mathcal A}=\Delta_0-c(p)\beta$.
\end{lemma}

\begin{proof}
$\partial_vF_{\mathcal A}=D_v$. By closedness $\partial_uD_v=\partial_vD_u$, so
$\partial_uF_{\mathcal A}=D_u(u,0)+\int_0^v\partial_tD_u(u,t)\,dt=D_u(u,v)$. Periodicity in $v$: $\int_0^1D_v(u,t)\,dt$ is the period
of $\Delta_0-c\beta$, which is $c(p)-c(p)\cdot1=0$ by Proposition~\ref{prop:cub-handle}(c). Smoothness and affine dependence
follow from those of $D_u$, $D_v$.
\end{proof}

\begin{construction}[The global primitive $\alpha_1$]\label{cons:fer-alpha1}
\lean{SymplecticPacking.Bundle.Cubic.alpha1pre,
SymplecticPacking.Bundle.Cubic.Ajet,
SymplecticPacking.Bundle.Cubic.Bjet,
SymplecticPacking.Bundle.Cubic.jetSum,
SymplecticPacking.Bundle.Cubic.alpha1}
\uses{def:fer-Delta, def:fer-FA, prop:fer-handle, M-eq:hdl-alpha1}
Define the base one-form, with $\zeta_{\mathcal A}F_{\mathcal A}$ extended by zero outside $\mathcal A^+$,
\[
 \alpha_{\rm pre}(p)=\alpha^0(p)-c(p)\beta-d_\Sigma\big(\zeta_{\mathcal A}F_{\mathcal A}(\cdot,p)\big),
\]
write its value at $b_j$ as $A_j(p)\,dx_j+B_j(p)\,dy_j$ ($A_j,B_j$ affine in $p$), and put
\begin{equation}\label{eq:fer-alpha1}
 \alpha_1(p)=\alpha_{\rm pre}(p)-\sum_{j=1}^3d_\Sigma\Big[\zeta_j\big(A_j(p)\,x_j+B_j(p)\,y_j\big)\Big].
\end{equation}
We use $\alpha_1$ as the toric one-form $\pi_N^*\alpha_1(p(\cdot))$ on $N$. For the data of \Msec{sec:hdl} this is the
one-form (\Mnum{eq:hdl-alpha1}), with the base point $(\frac16,0)$ for $F_{\mathcal A}$ and the cutoffs
$\zeta_{\mathcal A}$, $\zeta_j$ of Proposition~\ref{prop:fer-handle}.
\end{construction}

\begin{lemma}[Properties of $\alpha_1$]\label{lem:fer-alpha1}
\lean{SymplecticPacking.Bundle.Cubic.smooth_alpha1,
SymplecticPacking.Bundle.Cubic.alpha1_periodic,
SymplecticPacking.Bundle.Cubic.alpha1_affine,
SymplecticPacking.Bundle.Cubic.alpha1_curvature,
SymplecticPacking.Bundle.Cubic.Gamma1,
SymplecticPacking.Bundle.Cubic.alpha1_annulus}
\uses{cons:fer-alpha1, lem:fer-FA, lem:fer-frame, prop:fer-handle}
\begin{enumerate}[label=(\roman*)]
\item $\alpha_1$ is a smooth one-form on $N$, horizontal (no $dp$ and no fibre-angle component), affine in $p$, and
$d_\Sigma\alpha_1=\kappa_1-\kappa$.
\item Over $\mathcal A$, in the frame $(Z_0/|Z_0|)^3$,
\begin{equation}\label{eq:fer-Gamma1}
 \Gamma_{\mathcal A}+\alpha_1=\Gamma_{\rm sm}+f(s)H^*(p)\beta,\qquad\text{hence}\qquad
 \Omega_0+d\alpha_1=\omega_V+e_h\nu+d\big(f(s)H^*(p)\beta\big),
\end{equation}
with the full exterior derivative, which contains the mixed term $f(s)H^{*\prime}(p)\,dp\wedge\beta=-9f(s)\,dp\wedge\beta$.
\end{enumerate}
\end{lemma}

\begin{proof}
(i) Each term is a base one-form whose coefficients are smooth functions of the base point and of $p$, affine in $p$; as $p$
is smooth on $N$, the pull-back is a smooth horizontal one-form. The terms $c(p)\beta$, $d_\Sigma(\zeta_{\mathcal A}F_{\mathcal A})$ and
the corrections at the $b_j$ are $d_\Sigma$-closed, so $d_\Sigma\alpha_1=d_\Sigma\alpha^0=\kappa_1-\kappa$.
(ii) On $\mathcal A$, $\zeta_{\mathcal A}=1$ on a neighbourhood and all $\zeta_j=0$ (Proposition~\ref{prop:fer-handle}), so
$\alpha_1=\alpha^0-c\beta-(\Delta_0-c\beta)=\alpha^0-\Delta_0$, which is the first identity by the definition of $\Delta_0$. By
Lemma~\ref{lem:fer-frame}, $\Omega_0+d\alpha_1=\omega_V+d(\Gamma_{\mathcal A}+\alpha_1)=\omega_V+d\Gamma_{\rm sm}+d(fH^*\beta)$, and
$d\Gamma_{\rm sm}=e_h\nu$ because $\Gamma_{\rm sm}$ does not depend on $p$.
\end{proof}

\begin{remark}[The period is subtracted in $\alpha_1$]\label{rem:fer-period}
The term $-c(p)\beta$ is part of $\alpha_1$, not only of $F_{\mathcal A}$: otherwise the left side of \eqref{eq:fer-Gamma1} gains
$c(p)\beta$, whose ambient differential $c'(p)\,dp\wedge\beta$ is in general nonzero. Since $\beta$ is closed on the torus and
$\operatorname{supp}\beta\cap U_j=\varnothing$, the subtraction changes neither the base curvature nor the values at the $b_j$; it
only changes the mixed term, which does not enter \eqref{eq:fer-positivity}.
\end{remark}

\begin{lemma}[Vanishing on three whole fibres]\label{lem:fer-alpha1-fibres}
\lean{SymplecticPacking.Bundle.Cubic.alpha1_fibre_zero}
\uses{cons:fer-alpha1, prop:fer-handle}
At every point of each of the three fibres $N_{b_1},N_{b_2},N_{b_3}$, $\alpha_1$ is the zero covector of $T^*N$.
\end{lemma}

\begin{proof}
On the disc of radius $\frac1{200}$ about $b_j$ we have $\zeta_j=1$, and the cutoffs $\zeta_i$, $i\ne j$, vanish on $U_j$
(Proposition~\ref{prop:fer-handle}). Hence near $b_j$ the $j$-th correction equals $d_\Sigma(A_j(p)x_j+B_j(p)y_j)=A_j(p)\,dx_j+B_j(p)\,dy_j$,
and the base covector $\alpha_1(p)_{b_j}=\alpha_{\rm pre}(p)_{b_j}-A_j(p)\,dx_j-B_j(p)\,dy_j$ vanishes for every $p$. At a point of
$N_{b_j}$ with moment $p_0$, the toric form $\alpha_1$ is the pull-back of $\alpha_1(p_0)_{b_j}$ by $d\pi_N$, hence zero.
\end{proof}

\begin{remark}[Ambient covector versus pull-back]\label{rem:fer-covector}
The pull-back of $\alpha_1$ to a fibre vanishes automatically, since $\alpha_1$ is horizontal. That weaker statement would only make
the Moser field $Y$ of the path ($\iota_Y\Omega_\lambda=-\alpha_1$) $\Omega_\lambda$-orthogonal to the fibre, and such a field is in general
transverse to it. Lemma~\ref{lem:fer-alpha1-fibres} gives $Y=0$ at every point of $N_{b_j}$, so the Moser flow fixes these fibres
pointwise and preserves their complement, which the $\Z/3$-equivariant transfer maps into $\{Z_0\ne0\}$
(Lemma~\ref{lem:fer-alpha1-fibres} and Proposition~\ref{prop:rtc-avoid}).
\end{remark}

%----------------------------------------------------------------------
\subsection{Positivity including the mixed terms}\label{sec:fer-positivity}

The corrections above change $\alpha_1$ by base forms whose coefficients depend on $p$, so their ambient differentials contain
mixed terms $dp\wedge(\cdots)$ that do not appear in the base curvature. They do not affect nondegeneracy: locally the space is a
two-dimensional base times a fibre, and a mixed term $dp\wedge a$ with $a$ a base one-form wedges to zero with every base
two-form (three base factors), with $\omega_V=dp\wedge d\theta$ (two factors $dp$), and with itself.

\begin{lemma}[Positivity]\label{lem:fer-positivity}
\lean{SymplecticPacking.Bundle.Cubic.OmegaT,
SymplecticPacking.Bundle.Cubic.kappaT,
SymplecticPacking.Bundle.Cubic.OmegaT_split,
SymplecticPacking.Bundle.Cubic.OmegaT_sq,
SymplecticPacking.Bundle.Cubic.OmegaT_nondeg',
SymplecticPacking.Bundle.Cubic.OmegaT_nondeg'}
\uses{lem:fer-alpha1, lem:fer-kappa, lem:fer-Omega0, lem:pre-top-power}
For $\lambda\in[0,1]$ let $\Omega_\lambda=\Omega_0+\lambda\,d\alpha_1$ and $\kappa_\lambda=(1-\lambda)\kappa+\lambda\kappa_1$. In every local unitary
frame,
\[
 \Omega_\lambda=\omega_V+\kappa_\lambda+dp\wedge a_\lambda,\qquad a_\lambda=a_{\rm loc}+\lambda\,\partial_p\alpha_1,
\]
with $a_\lambda$ a base one-form, and
\begin{equation}\label{eq:fer-positivity}
 \Omega_\lambda^2=2\,\omega_V\wedge\kappa_\lambda ,
\end{equation}
smoothly across the zero section. Hence $\Omega_\lambda$ is a circle-invariant symplectic form on the open set
$\{p<\frac13,\ H^*(p)>0\}=\{p<\frac13-\frac{e_h}9\}$, which contains $N_{p\le h}$; moreover $\iota_{X_p}\Omega_\lambda=-dp$ for the
generator $X_p$ of the circle action, and $\partial_\lambda\Omega_\lambda=d\alpha_1$.
\end{lemma}

\begin{proof}
By Lemma~\ref{lem:fer-alpha1}(i), $d\alpha_1=d_\Sigma\alpha_1+dp\wedge\partial_p\alpha_1=\kappa_1-\kappa+dp\wedge\partial_p\alpha_1$, and with
\eqref{eq:fer-Omega0-local} this gives the local form. By the paragraph preceding the lemma (Lemma~\ref{lem:pre-top-power}(a) with
$m=1$), $\Omega_\lambda^2=2\,\omega_V\wedge\kappa_\lambda$. Write $\kappa_\lambda=k_\lambda\nu$. Where $p<\frac13$ and $H^*(p)\ge0$, $k_\lambda>0$ by
Lemma~\ref{lem:fer-kappa}, and $\omega_V\wedge\nu$ is a positive volume form; $H^*(p)>0$ iff $p<\frac13-\frac{e_h}9$, and
$h<\frac13-\frac{e_h}9$ because $e_h\le\frac32(1-3h)<9(\frac13-h)$. In Cartesian fibre coordinates $\omega_V=d\xi\wedge d\eta$ and
$dp=2\pi(\xi\,d\xi+\eta\,d\eta)$ are smooth, so the decomposition and \eqref{eq:fer-positivity} hold across $w=0$. The forms $\Omega_0$
and $\alpha_1$ are invariant under the circle action; $\iota_{X_p}\Omega_0=-dp$ (Lemma~\ref{lem:pre-local-form}(c)) and
$\iota_{X_p}d\alpha_1=\mathcal L_{X_p}\alpha_1-d(\alpha_1(X_p))=0$ because $\alpha_1$ is invariant and horizontal.
\end{proof}

Lemmas~\ref{lem:fer-alpha1}, \ref{lem:fer-alpha1-fibres} and~\ref{lem:fer-positivity} give the first path and the full marked-fibre vanishing for the cubic model; \eqref{eq:fer-positivity} is the four-dimensional case of the toric
top-power identity. The twist, the second and third paths, their primitives and the Moser map are constructed in
\S\ref{sec:shd-four}; all three variation primitives vanish on the marked fibres (Lemma~\ref{lem:shd4-fibres}), so
the Moser map fixes them pointwise (Proposition~\ref{prop:shd4-moser}).


\begin{remark}[Scalar data for both four-dimensional branches]\label{rem:fer-lean}
\lean{SymplecticPacking.Branch2.CubicData.Valid.Hstar_ge,SymplecticPacking.Branch2.CubicData.Valid.residual_ge}
The equal branch uses $g=3(1-t_0)$ and $h=t_0/3$, so $g=\Hb(h)$; the general branch uses $g=3d/4$ with the scalar
comparison datum of Proposition~\ref{prop:scalar-cubic-data}. In both cases the data are valid in the sense of
Definition~\ref{def:shd4-data}: $g\le\Hb(h)$, and \eqref{eq:shd4-comparison} holds with constant $\frac{31}{32}g$. With
$e_h=g/16$ and the cap smoothing scale $\mu_h=g/(16k)$, the smoothing loss is at most $g/32$ and the residual gap is
\[
 \frac{31g}{32}-\frac g{16}-\frac g{32}=\frac{7g}{8}>0.
\]
This is Lemma~\ref{lem:shd4-budget}(a),(c). The cover primitives, the marked-fibre vanishing and the positivity proof
above therefore apply to both branches.
\end{remark}
